% Part: incompleteness
% Chapter: theories-computability
% Section: intepretability

\documentclass[../../../include/open-logic-section]{subfiles}

\begin{document}

\olfileid[fr]{inc}{tcp}{itp}

\olsection{Les théories dans lesquelles $\Th{Q}$ est interprétable sont indécidables}

Nous pouvons renforcer encore ces résultats.
Informellement, interpréter un langage~$\Lang{L_1}$
dans un langage~$\Lang{L_2}$ consiste à définir,
au moyen de !!{formula}s de~$\Lang{L_2}$,
un domaine non vide ainsi que les symboles de relation et
de fonction de~$\Lang{L_1}$, par une traduction calculable qui préserve la déduction.
Les conditions assurant que les fonctions sont bien définies font partie de l'interprétation; nous ne la formaliserons pas ici.

\begin{thm}
Supposons que $\Th{T}$ soit une théorie dans un
langage où l'on peut interpréter le langage de
l'arithmétique de telle sorte que $\Th{T}$ soit
cohérente avec l'interprétation de~$\Th{Q}$.
Alors $\Th{T}$ est indécidable. Si $\Th{T}$
démontre l'interprétation des axiomes de~$\Th{Q}$,
aucune extension cohérente de~$\Th{T}$ n'est
décidable.
\end{thm}

\begin{proof}
Notons $!A^I$ la traduction d'un énoncé arithmétique~$!A$,
et soit $!E$ la conjonction des huit axiomes de~$\Th{Q}$.
Si $\Th{T}$ était décidable, on pourrait décider
l'appartenance à l'ensemble
\[
C=\Setabs{!A}{\Th{T}\Proves !E^I\lif !A^I}
\]
des énoncés arithmétiques qui satisfont cette condition.
En effet, la reconnaissance des énoncés et la traduction
sont calculables.
Si $\Th{Q}\Proves !A$, le théorème de la déduction et
l'interprétation donnent $\Th{T}\Proves !E^I\lif !A^I$.
Si $\Th{Q}\Proves\lnot !A$, ils donnent
$\Th{T}\Proves !E^I\lif\lnot !A^I$. Dans ce second cas,
$!A$ ne peut appartenir à~$C$, car $\Th{T}$ serait alors
incohérente avec $!E^I$. Ainsi, $C$ séparerait les
énoncés démontrables dans~$\Th{Q}$ de ses énoncés
réfutables, contrairement au lemme d'inséparabilité.
Si les axiomes interprétés sont déjà démontrables dans~$\Th{T}$,
toute extension cohérente de~$\Th{T}$ reste cohérente
avec eux. Le premier résultat s'applique donc à cette extension.
\end{proof}
La théorie des ensembles de Zermelo--Fraenkel avec
l'axiome du choix, $\Th{ZFC}$, peut servir de
fondement axiomatique assez puissant pour développer
une grande partie des mathématiques ordinaires.
On peut y définir les entiers naturels et y démontrer,
sous cette interprétation, les axiomes de~$\Th{Q}$.
Nous obtenons donc les corollaires suivants.

\begin{cor}
Aucune extension cohérente de~$\Th{ZFC}$ n'est
décidable.
\end{cor}

\begin{cor}
Il n'existe aucune extension de~$\Th{ZFC}$ qui soit
complète, cohérente et admette une axiomatisation
calculable.
\end{cor}

Le langage de~$\Th{ZFC}$ ne contient qu'un symbole
de relation binaire, $\in$; l'égalité elle-même
n'est pas nécessaire. Un résultat de codage des langages donne aussi :\footnote{L'indécidabilité de la logique pure s'établit par réduction de langages, indépendamment de la cohérence de~$\Th{ZFC}$.}

\begin{cor}
La logique du premier ordre de tout langage qui
possède un symbole de relation binaire est
indécidable.
\end{cor}

Ce résultat s'étend aux langages avec égalité et deux symboles
de fonction unaires, qui permettent de simuler un
symbole de relation binaire. Les résultats précédents
sont précis quant à la frontière : la logique du
premier ordre, avec égalité, d'un langage dont les seuls symboles non logiques sont
des relations \emph{unaires}, des constantes et au plus un symbole de
fonction \emph{unaire} est décidable.

Encore une remarque. Nous savons que l'ensemble des
!!{sentence}s vrais dans le modèle standard, dans
le langage $\Obj 0$, $'$, $+$, $\times$, $<$,
est indécidable. On peut définir $<$ au moyen des
autres symboles, puis $+$ au moyen de $\times$ et
de~$'$. L'ensemble des !!{sentence}s vrais dans
le langage $\Obj 0$, $'$, $\times$ est donc
indécidable. En revanche, Presburger a montré que
l'ensemble des !!{sentence}s du langage $\Obj 0$,
$'$, $+$ vrais dans l'arithmétique standard est
décidable, bien que le coût du calcul soit très élevé
dans le pire des cas.

\begin{explain}
La cohérence exclut le cas où tous les énoncés sont
démontrables. Pour une extension cohérente, les axiomes
arithmétiques interprétés restent démontrables, et le résultat
d'indécidabilité s'applique à cette extension.
\end{explain}

\end{document}
